% !TEX root = userguide.tex

\def\headdate{31th December 2008}

\section{Fluid-solid interaction}
\label{sec:fluidsolid}

In this section we will briefly describe fluid-solid interaction and give some
examples.

\subsection{Interaction constraint}

In fluid-solid interaction problems we have a fluid region $\Omega_\text{f}$
and a solid region $\Omega_\text{s}$. The interface between the fluid and the
solid is denoted by $\Gamma_\text{sf}$.
\begin{center}
   \includegraphics[scale=0.5]{fluid-solid}
\end{center}

The fluid will be described by the velocity $\vek u(\vek x,t)$ and 
additional quantities such as pressure $p(\vek x,t)$ for incompressible fluids
and structure parameters in case of complex fluids.
Here we assume that the solid is elastic
and can be described by the displacement $\vek d(\vek x,t)$ with respect to
some reference configuration. If the solid is incompressible we need the 
quantity pressure $p(\vek x,t)$ as well.

Both the fluid and solid region is governed by the mass balance, 
momentum balance, constitutive equation and suitable boundary and initial
conditions. In addition we assume the `no-slip' condition and a force balance
at the interface $\Gamma_\text{sf}$ between fluid and solid:
\begin{align}
   \vek t_\text{f}+\vek t_\text{s} &= \vek 0 \label{eq:forcebalance} \\
    \vek u - \dot{\vek d} &= \vek 0 \label{eq:noslip}
\end{align}
where $\vek t_\text{f}$ is the traction on the fluid,
$\vek t_\text{s}$ the traction on the solid,
$\vek u$ is the velocity of the fluid and $\vek d$ is the displacement of
the solid relative to the reference configuration.

We will impose the interface conditions Eqs.~\eqref{eq:forcebalance} and
\eqref{eq:noslip} using a Lagrange multiplier. However to avoid derivatives
in constraints we first apply time-discretization 
to Eq.~\eqref{eq:noslip}, for example:
\begin{align}
\intertext{Euler implicit}
   \vek u_{n+1}-\frac{\vek d_{n+1}-\vek d_n}{\Delta t}&=\vek 0
   \label{eq:constr_solid_fluid_Euler}\\
\intertext{Crank-Nicolson}
   \frac12(\vek u_{n+1}+\vek u_n)-\frac{\vek d_{n+1}-\vek d_n}{\Delta t}&=\vek 0
   \label{eq:constr_solid_fluid_CN}\\
\intertext{second-order backwards differencing (Gear)}
   \vek u_{n+1}-\frac{\frac32\vek d_{
              n+1}-2\vek d_n+\frac12\vek d_{n-1}}{\Delta t}&=\vek 0
   \label{eq:constr_solid_fluid_BDF2}
\end{align}
where $\Delta t=t_{n+1}-t_n$ is the time step. Note, that the time-derivative
of $\vek d$ is in a Lagrangian sense and that the corresponding velocity
at the interface needs to be evaluated in that sense.
Therefore the Gear method is easier to implement than the Crank-Nicolson method
in a fictitious domain setting if second-order time stepping is needed.

Focusing now on the Euler implicit scheme
we have the following constraint term in the
saddle point problem:
\begin{equation}
   \vek\lambda\cdot(\vek u_{n+1}-\frac{\vek d_{n+1}-\vek d_n}{\Delta t})
\label{eq:saddlepoint}
\end{equation}
where $\vek\lambda$ is a Lagrange multiplier. Variation leads to
\begin{equation}
   \delta\vek\lambda\cdot(\vek u_{n+1}-\frac{\vek d_{n+1}-\vek d_n}{\Delta t})+
   \delta\vek u_{n+1}\cdot\vek\lambda
   -\delta\vek d_{n+1}\cdot\frac{1}{\Delta t}\vek\lambda
\end{equation}
where we have ignored the dependence of $\vek u_{n+1}$ on $\vek d_{n+1}$ for
fictitious domain method\footnote{
Strictly speaking we need to take into account
the dependence of $\vek u_{n+1}$ on $\vek d_{n+1}$ for
the fictitious domain method. Variation of Eq.~\eqref{eq:saddlepoint}
leads to
\[
   \delta\vek\lambda\cdot(\vek u_{n+1}-\frac{\vek d_{n+1}-\vek d_n}{\Delta t})+
   \delta\vek u_{n+1}\cdot\vek\lambda-
   \delta\vek d_{n+1}\cdot\big(-(\nabla\vek u)_{n+1}^T+
               \frac{1}{\Delta t}\ten{I}\big)\cdot\vek\lambda
\]
The additional term involving the velocity gradient gives an additional forcing
on the solid. However compared with the other term in the solid forcing, it will
vanish for $\Delta t\rightarrow0$. Therefore we think this term might be
an artifact of putting a constraint on a time discretized equation.
}.
The forcing on the fluid will be $\vek\lambda$, whereas the forcing on
the solid will be $-\vek\lambda/\Delta t$. Therefore we need
to multiply the momentum balance of the solid with the factor $1/\Delta t$ to
get the correct force balance Eq.~\eqref{eq:forcebalance} at the fluid-solid
interface. Note that this factor depends on the time-discretization used!

We will apply 
an iteration scheme to require equilibrium for the solid at each time step.
The constraint, like Eq.~\eqref{eq:constr_solid_fluid_Euler}, is in fact a
nonlinear equation if the fictitious domain method is applied,
since $\vek u_{n+1}$ depends on $\vek d_{n+1}$ in that case. However, we will
not take this into account in the linearization and therefore we will apply a
Picard like iteration scheme for the constraint.

\subsection{A fluid flowing in a channel with an elastic beam attached to the
wall: fluid-solid interaction}
\label{sec:elastic2}
                                                                                
We consider the Stokes equation for a 2D channel problem. Attached to the wall
there is an elastic beam that bends due to the friction forces of the fluid
imposed on the solid. The flow is periodically reversed.
The interaction between the fluid and the solid is
modeled using a constraint on a `two-sided' object (see
Sec.~\ref{sec:objects}). The constraint is put on the surface of the solid.

The object consists of
points that are positioned at the nodal points of the solid\footnote{The points
on the object does not have to be equal to the nodal points of the solid, but
that is what we did here.}
The constraint equation
Eq.~\eqref{eq:noslip} will be discretized using collocation
and solved in a Lagrangian way. The
points on the object move with the solid and refer to the
same material point at any time. Eq.~\eqref{eq:noslip} is
discretized using an implicit Euler scheme as given by
Eq.~\eqref{eq:constr_solid_fluid_Euler}.
At each time step there is
an iteration scheme to require equilibrium at the end of the time step. The
iteration scheme uses the linear viscous matrix, the Newton-Raphson
linearization matrix of the solid and a Picard scheme for the constraint
equation Eq.~\eqref{eq:constr_solid_fluid_Euler}.

The example file is \texttt{elastic2.f90} and can be obtained by typing at the
command line:
\begin{quote}
   \texttt{getexample elastic}
\end{quote}
The input file is \texttt{input\_elastic2.txt}.
                                                                                
Some notes on the source:
\begin{listing}{396}

!   use boundary curve

    call add_to_mesh ( mesh, object='curve', objectcurve=9, typeofobject=2, &
      excludegroups=[2], excludegroups2=[1], excludecurves=[5] ) 

\end{listing}
Here the `two-sided' objects is created that connects the fluid with the solid.
\begin{listing}{428}

! create updatemesh from real mesh:
!   mesh: the original mesh. For the elastic solid the reference configuration.
!   updatemesh: copy of the original mesh, except for the
!               coordinates of the nodal points, which are redefined and
!               updated after each time step.
  call copy ( mesh, updatemesh )
  call fill_mesh_parts ( updatemesh )

\end{listing}
Here a mesh \texttt{updatemesh} is created by copying all data in the
original mesh. The coordinates of \texttt{updatemesh} will be
updated after each time step so that it can be in plotting.
Note, that \texttt{copy} is a generic routine that in this case copies 
\texttt{mesh} to \texttt{updatemesh}, but only the `basic' components.
Therefore, \texttt{fill\_mesh\_parts} is necessary to make a full mesh.
\begin{listing}{464}
! constraint between the fluid and the solid

  call define_constraint ( mesh, input_probdef, object=1, &
    discretization='collocation', nodedof=2, physq=1 )

\end{listing}
Here the constraint is defined on the `two-sided' objects.
\begin{listing}{476}

! all displacements of the solid
  call create_subscript ( mesh, problem, disp, physqarr=[1], groups=[2], &
    fillnodes=.true. )
\end{listing}
Here a vector subscript is created that points to the displacement degrees of
freedom in the solid. In the subscript (\texttt{disp}) also the nodes involved
in the subscript are stored to be used later in updating the coordinates in 
\texttt{updatemesh}.
\begin{listing}{555}

!     fluid
      call build_system ( mesh, problem, sysmatrix, rhsd, elemsub=stokes_elem, &
        coefficients=coefficients_f, elgroup1=1 )

!     solid
      call build_system ( mesh, problem, sysmatrix, rhsd, &
        elemsub=elastic_elem, coefficients=coefficients_s, &
        oldvectors=oldvectors, elgroup1=2, addmatvec=.true., &
        factormat=1._dp/deltat, factorvec=1._dp/deltat )

!     coupling
      call build_system_constraint ( mesh, problem, sysmatrix, rhsd, &
        elemsub=elementc, addmatvec=.true. )

\end{listing}
Here the system is build for the fluid, solid and constraint. Note, that the 
solid system is multiplied by a factor of $1/\Delta t$ to create a correct
force balance between fluid and solid.
\begin{listing}{613}

!   update coordinates of mesh

    updatemesh%coor(disp%nodes,:) = mesh%coor(disp%nodes,:) + &
           transpose ( reshape (soln%u(disp%s),[2,nnd]) )

\end{listing}
Here the coordinates in \texttt{updatemesh} are updated. Note the use of
\texttt{disp\%nodes}.

After running the program by
\begin{quote}
   \texttt{elastic2 < input\_elastic2.txt}
\end{quote}
we get the results shown in Figs.~\ref{fig:velocityfluid} and
\ref{fig:pressurebeam}.
Note, that the beam is `wet': there is also fluid inside the beam having the
same viscosity as the outside fluid. This creates a kind of viscoelastic solid.
Depending on the flexibility of the beam
the fluid stresses inside the beam can become important. These can be reduced
by lowering the viscosity inside the beam.
\begin{figure}[ht!]
   \begin{center}
   \includegraphics[scale=0.7]{velocity0200}
   \end{center}
   \caption{Velocity of the fluid and the deformed beam at time $t=10$.}
    \label{fig:velocityfluid}
\end{figure}
\begin{figure}[ht!]
   \begin{center}
   \hspace*{1.6cm}\includegraphics[scale=0.7]{pressure2_contour}
   \end{center}
   \caption{Pressure in the beam at time $t=10$.}
    \label{fig:pressurebeam}
\end{figure}

\subsection{A fluid flowing in a channel with an elastic beam attached to the
wall: fluid-solid interaction; weak constraints}
\label{sec:elastic4}

(These examples have been developed by Ioannis Dimakopoulos of the TU/e).

The examples are similar to the example in Sec.~\ref{sec:elastic4}, but now
it is possible to use a weak constraint for the fluid-solid interaction.
The example file \texttt{elastic4.f90} use constraints on the solid-fluid interface, whereas
the example file  \texttt{elastic5.f90} uses
 a constraint distributed on the whole solid domain and a velocity profile
sinusoidal in time.
